\documentclass[%
sor,
 jor,
 amsmath,amssymb,
 reprint,
longbibliography,
]{revtex4-2}
\usepackage{float}
\usepackage{caption}
\usepackage{empheq, amssymb, amsfonts}
\usepackage{siunitx}
\usepackage{graphicx}
\usepackage{circuitikz}
\usepackage{url}
\usepackage{hyperref}


\begin{document}
\title{PH3244\\Experiment - 7\\Frank Hertz}
\author{Aumshree P. Shah\\20231059}
\altaffiliation{\color{red}aumshree.pinkalbenshah@students.iiserpune.ac.in}
\date{\today}
\begin{abstract}
\centering
In this experiment we find the first ionizaion energy of a element in the Frank-Hertz tube.
\end{abstract}
\maketitle
\tableofcontents
\vspace*{\fill}
\pagebreak
\section{Theory and Procedure}
\subsection{Theory}




\subsection{Procedure}Refer to the manual [\cite{ph3244repo}] for a look at procedure, although the required pages will be attached at the back for sumbition along with the signed report.
\section{Observation}
\subsubsection{General Values\footnote{The least count is to be interpreted from the significant figures written}}
\begin{align*}
	\text{Model} &: \text{FH-300H}\\
	V_\text{G1K} &= 1.50 \si{\volt}\\
	V_\text{G2A} &= 7.50 \si{\volt}\\
\end{align*}
\subsection{Model 1 test data}
\begin{table}[H]
\begin{minipage}[t]{0.48\textwidth}
\centering
\begin{tabular}{|c|c|}
\hline
$V_{ork}$ (V) & $I\,(10^{-10}\,\mathrm{A})$ \\
\hline
2.0  & 0   \\
11.8 & 1   \\
13.6 & 3   \\
14.7 & 5   \\
16.9 & 12  \\
18.6 & 16  \\
19.5 & 18  \\
20.5 & 19  \\
22.7 & 20  \\
24.6 & 18  \\
25.6 & 19  \\
26.4 & 20  \\
27.3 & 25  \\
28.1 & 30  \\
29.4 & 40  \\
30.3 & 46  \\
31.5 & 50  \\
32.2 & 51  \\
32.8 & 50  \\
33.3 & 49  \\
34.0 & 47  \\
34.8 & 42  \\
35.7 & 35  \\
\end{tabular}
\end{minipage}%
\hfill % Adds horizontal stretchable space between minipages
\begin{minipage}[t]{0.48\textwidth}
\centering
\begin{tabular}{|c|c|}
36.3 & 33  \\
37.0 & 32  \\
37.8 & 38  \\
38.4 & 48  \\
38.9 & 58  \\
39.3 & 69  \\
39.6 & 78  \\
40.0 & 88  \\
40.8 & 107 \\
41.2 & 116 \\
41.7 & 123 \\
42.7 & 132 \\
43.1 & 132 \\
43.5 & 130 \\
44.3 & 119 \\
45.1 & 103 \\
45.7 & 88  \\
46.6 & 69-96\footnote{Here we seee that after some time even if the $V_\text{G2K}$ is constant the current inceases over time, hence in the next data we shall take it instantiously and will not allow any other effects to affect $I$.}\\
\hline
$V_{ork}$ (V) & $I\,(10^{-10}\,\mathrm{A})$ \\
\hline
\end{tabular}
\end{minipage}
\caption{Test data to see if there are any systematic errors}
\end{table}


\subsection{Model 2 instantious data}
\begin{table}[H]
\begin{minipage}[t]{0.48\textwidth}
\centering
\begin{tabular}{|c|c|}
\hline
$V_{ork}$ (V) & $I\,(10^{-10}\,\mathrm{A})$ \\
\hline
2.0 & 1 \\
10.7 & 2 \\
12.2 & 6 \\
14.0 & 11 \\
14.8 & 13 \\
16.4 & 13 \\
17.2 & 12 \\
18.3 & 10 \\
19.6 & 8 \\
20.1 & 7 \\
21.5 & 6 \\
23.0 & 8 \\
23.7 & 12 \\
24.7 & 17 \\
25.5 & 20 \\
26.3 & 22 \\
27.6 & 22 \\
28.8 & 20 \\
29.6 & 18 \\
30.6 & 15 \\
31.5 & 12 \\
32.4 & 8 \\
33.5 & 12 \\
35.7 & 18 \\
36.6 & 22 \\
37.4 & 26 \\
38.0 & 31 \\
39.0 & 32 \\
39.9 & 32 \\
40.8 & 29 \\
41.6 & 26 \\
42.3 & 23 \\
43.2 & 19 \\
44.1 & 17 \\
44.4 & 14 \\
46.8 & 16 \\
\end{tabular}
\end{minipage}%
\hfill % Adds horizontal stretchable space between minipages
\begin{minipage}[t]{0.48\textwidth}
\centering
\begin{tabular}{|c|c|}
47.8 & 21 \\
49.1 & 28 \\
50.5 & 37 \\
51.3 & 39 \\
52.3 & 40 \\
53.8 & 37 \\
55.2 & 31 \\
56.5 & 25 \\
57.8 & 21 \\
58.8 & 22 \\
59.8 & 25 \\
61.6 & 35 \\
62.6 & 40 \\
63.3 & 44 \\
64.1 & 47 \\
65.2 & 48 \\
66.2 & 48 \\
67.4 & 45 \\
68.2 & 42 \\
70.1 & 36 \\
71.2 & 35 \\
72.8 & 38 \\
73.9 & 43 \\
74.9 & 48 \\
76.4 & 55 \\
77.0 & 57 \\
79.0 & 61 \\
79.7 & 62 \\
80.4 & 62 \\ 
81.3 & 60 \\
82.8 & 58 \\
84.6 & 56 \\
85.0 & 58 \\
86.1 & 59 \\
87.3 & 63 \\
88.9 & 69 \\
90.8 & 78 \\
92.6 & 82 \\
92.8 & 83 \\
\hline
$V_{ork}$ (V) & $I\,(10^{-10}\,\mathrm{A})$ \\
\hline
\end{tabular}
\end{minipage}
\caption{Data taken instationously to reduce any systematic errors due to heating}
\end{table}
\begin{center}
	Data taken on 12 Mar 2026.
\end{center}

\section{Calculations, Analysis and Characteristics}
\subsection{Sources of Error}

\noindent % Prevents paragraph indentation before the first minipage
\begin{minipage}[t]{0.48\textwidth}
\subsubsection{Systematic Errors}
\begin{enumerate}
	\item In alignment of the beam with the crystal.
	\item Improper standing wave produced than the theoratical.
	\item Non Parallel Screen causing increase in distance.
	\item Index of fringe being differnet than what is measured. 
\end{enumerate}
\end{minipage}
\hfill % Distributes remaining horizontal space
\begin{minipage}[t]{0.48\textwidth}
\subsubsection{Random Errors}
\begin{enumerate}
	\item Human errros in measurement.
	\item Uncertanity in instruments.
	\item Shaking of table resulting in moving fringes.
	\item Diffraction increasing the fringe width.
\end{enumerate}
\end{minipage}

\subsection{Error Propogation}
To calculate the standard deviation from values we have measured with some least count, and the values we have applied digitally, we use the formula [\cite{tewari_uncertainty}]: 

\begin{equation}
\sigma = \frac{\text{Least Count}}{\sqrt{12}}
\end{equation}
and to calculate the standard deviation for values we have applied on analog instrument with some least count, we use the formula [\cite{tewari_uncertainty}]: 
\begin{equation}
\sigma = \frac{\text{Least Count}}{\sqrt{24}}
\end{equation}
Since all our observations only satisfy the criteria for Equation-6 we will be using that. Now from Equation-1,2,3 we can write $V$ as : 
\begin{equation}
	V = \mu n\lambda \frac{\sqrt{{D^2+L^2}}}{D}
\end{equation}
to calculate the error of $V$ we notice that it'll only (my only here I mean the most significant factor of error) will be the uncertainty in $D$, hence in error calculation we will assume other errors are negligible.
Now using formula for error propogaiton [\cite{taylor_error_analysis}], efor Equation-8 we see that: 
\begin{align*}
	\sigma_V / \sigma_D &=  \frac {\text d}{\text d D} V\\
		 &= \mu n \lambda \frac 1 D\frac {\text d}{\text d D}  \left( \sqrt{D^2+L^2} \right) + \mu n \lambda \sqrt{D^2+L^2}\frac {\text d}{\text d D}  \left( \frac 1 D \right) 
\end{align*}
\begin{equation}
	\implies \sigma_V	 =\boxed{\left( \mu n \lambda \frac {1}{\sqrt{D^2+L^2}} - \mu n \lambda \sqrt{D^2+L^2}\frac {1}{D^2} \right) \sigma_D}
\end{equation}
Similarly from Eqation-4,5 we get the uncertainty in $\beta$ and $K$ as:
\begin{equation}
	\sigma_\beta = 2V\rho \sigma_V
\end{equation}
\begin{equation}
	\sigma_K = \frac {2}{\rho V^3} \sigma_V
\end{equation}

\subsection{Calculations}

First we see that since the Table-2 has only 2 data points, where also there is a good chance that it might have high precision and low accuract, we wil not consider the data from Table-2 in our calculations. Using Equation-8,4,5 to calculate the variables and Equation-9,10,11 to calcualte their uncertainty, we get:
\begin{table}[H]
\centering
\begin{tabular}{c c c c c c c}
\hline
$n$ & Distance & Least Count & SI Distance & SI Uncertainty & Velocity & Velocity Uncertainty  \\
    & (cm) & (cm) & (m) & (m) & (m/sec) & (m/sec) \\
\hline
1 & 0.45 & 0.05 & 0.0045 & 0.00014 & 1343 & 42 \\
2 & 0.85 & 0.05 & 0.0085 & 0.00014 & 1422 & 23 \\
3 & 1.30 & 0.10 & 0.0130 & 0.00029 & 1394 & 29 \\
4 & 1.80 & 0.10 & 0.0180 & 0.00029 & 1343 & 19 \\
\hline
\end{tabular}

\caption{Calculating velocity from TABLE-1}
\end{table}
Using the weighted average formula [\cite{standard_error_wiki}] we can calculate $V$, $\beta$ and $K$ which will be discussed in results. 
\begin{equation}
	\bar{v}=\frac{\sum_{i=1}^{N} \frac{v_i}{\sigma_i^2}}{\sum_{i=1}^{N} \frac{1}{\sigma_i^2}}
\end{equation}
\begin{equation}
	\sigma_{\bar{v}} = \sqrt{\frac{1}{\sum_{i=1}^{N} \frac{1}{\sigma_i^2}}}
\end{equation}
e weighted average we get is: From this 

\section{Result and Conclusion}
By taking the weighted average of velocity and Equaiton 12,13 of each value from TABLE-3, we get value of $V$ as $\boxed{V = 1375 \pm 12 \text{ m/sec}}$; From Equation-4,10 we get the bulk modulus as $\boxed{\beta = 1.88 \pm 0.03 \text{ GPa}}$ and finally from Equation-5,11 we get the compressiblity as $\boxed{K = 53.67 \pm 0.08 \times 10^{-6}\text{ atm}}$ . We will compare this values with the theoratical values in Appendix-A.

\vspace{1cm}
\hrule
\vspace{1cm}
\appendix

\bibliographystyle{apsrev4-2}
\bibliography{refrences}
\hfill\pagebreak

\section{Theoratical values and comparision}
The theoratical value for bulk modulus of water at 20\si{\celsius} is around 2.2 GPa, this is 16\% different from the calculated while the uncertainty being only 1\%. The reason for this are the systematic errors discussed in the error analysis section.
\section{Improvments in procedure}
We can measure the $D$ on a flourosent paper in dark room which will make the readings a lot easier and accurte. Furthermore we can calculate the $D$ at very large $L$ this will make the relative uncertainty of $D$ small and thus giving more accurate readings. For the random error, we can take more readings to reduce them.
\end{document}
